\noindent\hspace*{\parindent}I give the following further explanations of the individual items in the preceding balance account.

For (1). The capital value of the fixed part of the widows' pensions for all 15 widows currently living (at 200 thalers for each) is the sum of the present capital values of these pensions, calculated separately for each widow according to her age, with allowance for possible orphan pensions where there are currently minor children. I give these values individually below: the widows are numbered according to the full sequence since the institution was founded.
\begin{center}
\footnotesize\setlength{\tabcolsep}{3pt}
\begin{tabular}{r|l|r l||r|l|r l||r|l|r l}
31&\textsc{Sch.}&1066&thalers&55&H.&1834&thalers&65&H.&2003&thalers\\
44&F.&1622&&58&G.&2730&&66&H.&1172&\\
50&H.&2849&&59&W.&1862&&68&M.&2923&\\
52&P.&1479&&62&B.&1003&&69&D.&2897&\\
53&M.&1812&&63&G.&1309&&70&L.&2224&
\end{tabular}
\end{center}
\noindent\hspace*{\parindent}The figures (2) and (7) were obtained by the following procedure, whose justification can be found in the memorandum of 1845. Of the 53 members who currently make up the association, 47 are married at present. For each of them, according to the ages of the husband and wife, the present capital value has been calculated both of the contributions to be paid by the husband (at 15 thalers annually) and of the fixed part of the widow's pension that the wife would receive if she survives him (at 200 thalers). Understood in this sense, the sum of the contributions comes to 7947 thalers, and the sum of the pensions to 45364 thalers. To take account of the 6 members who are currently unmarried, these figures are multiplied by \(\tfrac{53}{47}\), giving contributions \(=8961\) thalers and pensions \(=51155\) thalers: the former figure is item (7) above. The latter, increased by \(\tfrac{1}{7}\) on account of orphan pensions, gives 58463 thalers, item (2).

The justification for adopting the fraction \(\tfrac{1}{7}\) for orphan pensions, instead of the fraction used in 1845, \(\tfrac{1}{6}\), will be given when items (3) and (8) are explained.

I also give below the results of the above calculation for each of the 47 married members. The numbering follows the order of entry into the institution; the figures in the third column are the contributions, and those in the fourth are the pensions.

\SourceLeaf{171}
\begin{center}
\footnotesize\setlength{\tabcolsep}{2pt}
\begin{tabular}{r|l|r|r||r|l|r|r||r|l|r|r}
135&O.&102 Thl.&1179 Thl.&170&R.&124 Thl.&1261 Thl.&199&F.&167 Thl.&937 Thl.\\
138&C.&82&1402&171&C.&132&955&200&L.&211&748\\
139&U.&137&730&173&F.&168&931&201&W.&194&749\\
140&H.&134&1199&174&H.&163&1354&202&M.&201&991\\
141&L.&113&833&175&R.&195&996&203&B.&167&751\\
147&R.&133&1188&176&L.&190&1122&204&E.&201&697\\
149&G.&127&1439&177&W.&176&936&205&H.&195&790\\
151&O.&173&870&180&T.&177&887&206&H.&188&854\\
153&B.&154&1271&185&W.&198&919&207&W.&200&783\\
155&K.&145&727&186&H.&169&1134&208&E.&175&1278\\
156&v. S.&146&662&187&W.&180&748&209&B.&154&1051\\
160&H.&161&1173&189&H.&172&953&211&S.&195&1164\\
161&B.&156&682&192&R.&188&782&212&T.&181&895\\
164&W.&159&1155&193&G.&206&896&213&W.&200&811\\
167&S.&190&835&195&D.&191&822&214&B.&188&915\\
169&Z.&180&943&196&B.&207&968&&&&
\end{tabular}
\end{center}
